\chapter{补充图表与完整结果文件说明}

\section{补充图表}

图 \ref{fig:appendix-numeric-histograms} 展示主要数值变量分布，作为探索性分析的补充材料。

\maybefigure{eda_numeric_histograms.png}{主要数值变量分布直方图}{fig:appendix-numeric-histograms}

图 \ref{fig:appendix-behavior-scatter} 展示行为变量与结果变量之间的关系，用于辅助观察变量联系，不作为因果证据。

\maybefigure{eda_behavior_outcome_scatter.png}{行为变量与结果变量关系图}{fig:appendix-behavior-scatter}

图 \ref{fig:appendix-productivity-prediction} 展示 \texttt{productivity\_score} 的真实值与预测值对比。该图与正文中的负结果分析一致，说明生产力得分在当前特征下预测较弱。

\maybefigure{regression_productivity_observed_vs_predicted.png}{生产力得分真实值与预测值对比}{fig:appendix-productivity-prediction}

图 \ref{fig:appendix-classification-importance} 展示分类任务置换重要性，用于补充说明模型主要依赖的输入变量。

\maybefigure{classification_permutation_importance.png}{分类任务置换重要性}{fig:appendix-classification-importance}

\section{PCA 解释方差表}

表 \ref{tab:pca-explained-variance} 列出前 10 个主成分的解释方差比例。PCA 仅用于可视化和辅助理解，不作为分类或回归模型输入。

\begin{table}[H]
\centering
\caption{PCA 解释方差附录表（前 10 个主成分）}
\label{tab:pca-explained-variance}
\begin{tabular}{rrr}
\toprule
主成分 & 解释方差比例 & 累计解释方差 \\
\midrule
PC1 & 0.2694 & 0.2694 \\
PC2 & 0.1547 & 0.4241 \\
PC3 & 0.1366 & 0.5607 \\
PC4 & 0.1258 & 0.6865 \\
PC5 & 0.1030 & 0.7895 \\
PC6 & 0.0759 & 0.8654 \\
PC7 & 0.0430 & 0.9084 \\
PC8 & 0.0425 & 0.9509 \\
PC9 & 0.0274 & 0.9783 \\
PC10 & 0.0152 & 0.9935 \\
\bottomrule
\end{tabular}
\end{table}


\section{完整结果文件说明}

完整结果文件保存在项目根目录的 \texttt{results/} 中。正文只排版精简表格，完整 CSV 作为电子附录保留，包括：

\begin{itemize}
  \item \texttt{classification\_cv\_results.csv}：分类超参数搜索完整结果。
  \item \texttt{classification\_tuned\_metrics.csv}：分类阈值策略测试集指标。
  \item \texttt{regression\_digital\_dependence\_cv\_results.csv}：数字依赖回归调参结果。
  \item \texttt{regression\_productivity\_metrics.csv}：生产力弱预测结果。
  \item \texttt{clustering\_model\_comparison.csv}：聚类算法和簇数对比。
  \item \texttt{clustering\_lifestyle\_profiles.csv}：完整聚类画像表。
  \item \texttt{pca\_explained\_variance.csv}：完整 PCA 解释方差表。
\end{itemize}
